\section{Branching with variable zeroing}

% Motivation for variable zeroing
There are methods !!!Add citations!!! in the literature of algorithmic algebraic geometry that assume the value of a variable is zero, and then solve the resulting system.

% Explaining the method with variable zeroing
The following method takes the algorithm~\ref{alg:Branching} and allows for the solver to time out.
If the solver times out, the method chooses a variable, and tries to solve the system using algorithm~\ref{alg:Branching} once assuming the variable is zero, and once assuming the variable is non-zero.
The branching procedure is then applied recursively either until every branch is resolved, or until the computation resources are exhausted.
In the experiments section~\ref{sec:experiments}, we show that recursive branching~\ref{alg:ZeroingBranching} utilizes the computation resources better than the original branching method~\ref{alg:Branching}.

% Pseudocode for the method with variable zeroing
\begin{algorithm}
\footnotesize
\caption{Recursive branching with variable zeroing}
\label{alg:ZeroingBranching}
\begin{algorithmic}[1]
    \Require A polynomial system $P$, and a timeout $\tau$.
    \Function{ZeroingBranching}{$P, \tau$}
        \State $C \gets \emptyset$ \Comment{One global unsatisfiable-core set}
        \State $X \gets \Call{Variables}{P}$ \Comment{All variables of $P$}
        \State \Return $\Call{SolveByZeroing}{P,X,\tau,C}$
    \EndFunction
    \Function{SolveByZeroing}{$P,X,\tau,C$}
        \State $(S,T) \gets \Call{BranchingWithTimeout}{P,\tau,C}$
        \Comment{Algorithm~\ref{alg:Branching}; $C$ is shared and updated}
        \State $U \gets \emptyset$ \Comment{Branches left unresolved}

        \ForAll{$Q \in T$}
            \If{$X=\emptyset$}
                \State $U \gets U\cup\{Q\}$
                \State \textbf{continue}
            \EndIf

            \State $x \gets \Call{ChooseVariable}{Q,X}$
            \State $(S_0, U_0, C) \gets \Call{SolveByZeroing}{Q\cup\{x\}, X\setminus\{x\}, \tau,C}$
            \State $(S_1, U_1, C) \gets \Call{SolveByZeroing}{Q\cup\{1-xz\}, X\setminus\{x\}, \tau,C}$
            \Comment{$z$ is fresh}

            \State $S \gets S\cup S_0\cup S_1$
            \State $U \gets U\cup U_0\cup U_1$
        \EndFor

        \State \Return $(S, U, C)$
    \EndFunction
\end{algorithmic}
\end{algorithm}

Solutions of a $Q_1$ branch are projected onto the original variables, so the fresh variables do not appear in the output.

The run on \texttt{systems/system.txt}, with a five-second timeout for each
cover, produces the recursion tree in Figure~\ref{fig:variable_zeroing_tree}.
The variable chosen after a timeout is written inside the corresponding cover
node.  Its outgoing edges give the zero and nonzero cases.

\begin{figure}[htbp]
    \centering
    \resizebox{\textwidth}{!}{% Source: results/physical_system_m4gb_5s_figure_run.json
% Plot-only changes: omit unused core-pruning legend; label outcomes SAT/UNSAT.
% Input: system.txt
\providecommand{\vzInconsistent}{\ensuremath{\textcolor{red!75!black}{\boldsymbol{\times}}}}
\providecommand{\vzConsistent}{\ensuremath{\textcolor{green!45!black}{\boldsymbol{\checkmark}}}}
\providecommand{\vzNoCovers}{\ensuremath{\textcolor{gray!75!black}{\boldsymbol{\varnothing}}}}
\providecommand{\vzTimeout}{\tikz[baseline=-0.55ex,scale=0.105]{\filldraw[fill=yellow!70,draw=orange!85!black,line width=1.2pt] (0,0) circle (1);\draw[orange!85!black,line width=1.4pt,line cap=round] (0,0) -- (0,0.58) (0,0) -- (0.48,-0.28);}}
\providecommand{\vzContracted}{\ensuremath{\textcolor{violet!80!black}{\boldsymbol{\leadsto}}}}
\providecommand{\vzError}{\ensuremath{\textcolor{red!75!black}{\boldsymbol{!}}}}

\begin{tikzpicture}[
  font=\small,
  cover/.style={draw=blue!65!black,fill=blue!3,rounded corners=2pt,line width=0.65pt,minimum width=13mm,minimum height=7mm,inner sep=2pt,align=center},
  timeout cover/.style={cover,draw=orange!85!black,fill=yellow!12,minimum width=18mm,minimum height=10mm},
  contracted cover/.style={cover,draw=violet!75!black,fill=violet!5,minimum width=18mm,minimum height=10mm},
  error cover/.style={cover,draw=red!75!black,fill=red!5,minimum width=18mm,minimum height=10mm},
  empty branch/.style={draw=gray!65,dashed,fill=gray!5,rounded corners=2pt,minimum width=15mm,minimum height=7mm,inner sep=2pt,align=center},
  input/.style={draw=blue!70!black,fill=blue!8,circle,line width=0.7pt,minimum size=8mm},
  branch/.style={-{Stealth[length=1.8mm]},line width=0.65pt},
  stem/.style={line width=0.65pt},
  edge label/.style={fill=white,inner sep=1pt,font=\scriptsize}
]
\coordinate (group-0) at (0.000,-0.730);
\node[cover] (cover-0-0) at (-3.300,-1.450) {\vzInconsistent};
\node[cover] (cover-0-1) at (-1.100,-1.450) {\vzInconsistent};
\node[cover] (cover-0-2) at (1.100,-1.450) {\vzConsistent};
\node[timeout cover] (cover-0-3) at (3.300,-1.450) {\vzTimeout\\[-1pt]$\mathit{pph}$};
\coordinate (group-1) at (-1.500,-2.630);
\node[cover] (cover-1-0) at (-2.600,-3.350) {\vzInconsistent};
\node[timeout cover] (cover-1-1) at (-0.400,-3.350) {\vzTimeout\\[-1pt]$\mathit{DpDph2}$};
\coordinate (group-2) at (-7.800,-4.530);
\node[cover] (cover-2-0) at (-10.000,-5.250) {\vzInconsistent};
\node[cover] (cover-2-1) at (-7.800,-5.250) {\vzInconsistent};
\node[cover] (cover-2-2) at (-5.600,-5.250) {\vzInconsistent};
\coordinate (group-3) at (-2.600,-4.530);
\node[timeout cover] (cover-3-0) at (-2.600,-5.250) {\vzTimeout\\[-1pt]$\mathit{DpDph1}$};
\coordinate (group-4) at (-7.800,-6.430);
\node[empty branch] (pending-4) at (-7.800,-7.150) {\vzNoCovers\ no covers};
\coordinate (group-5) at (-4.800,-6.430);
\node[contracted cover] (cover-5-0) at (-4.800,-7.150) {\vzContracted\\[-1pt]$\mathit{xr}$};
\coordinate (group-6) at (-4.500,-8.330);
\node[cover] (cover-6-0) at (-5.600,-9.050) {\vzInconsistent};
\node[cover] (cover-6-1) at (-3.400,-9.050) {\vzInconsistent};
\coordinate (group-7) at (-0.400,-8.330);
\node[cover] (cover-7-0) at (-0.400,-9.050) {\vzInconsistent};
\coordinate (group-8) at (2.600,-2.630);
\node[timeout cover] (cover-8-0) at (2.600,-3.350) {\vzTimeout\\[-1pt]$\mathit{DpDph2}$};
\coordinate (group-9) at (3.700,-4.530);
\node[cover] (cover-9-0) at (0.400,-5.250) {\vzInconsistent};
\node[cover] (cover-9-1) at (2.600,-5.250) {\vzInconsistent};
\node[cover] (cover-9-2) at (4.800,-5.250) {\vzInconsistent};
\node[cover] (cover-9-3) at (7.000,-5.250) {\vzInconsistent};
\coordinate (group-10) at (10.000,-4.530);
\node[timeout cover] (cover-10-0) at (10.000,-5.250) {\vzTimeout\\[-1pt]$\mathit{xr}$};
\coordinate (group-11) at (1.500,-6.430);
\node[cover] (cover-11-0) at (-1.800,-7.150) {\vzInconsistent};
\node[cover] (cover-11-1) at (0.400,-7.150) {\vzInconsistent};
\node[cover] (cover-11-2) at (2.600,-7.150) {\vzInconsistent};
\node[cover] (cover-11-3) at (4.800,-7.150) {\vzInconsistent};
\coordinate (group-12) at (7.800,-6.430);
\node[timeout cover] (cover-12-0) at (7.800,-7.150) {\vzTimeout\\[-1pt]$\mathit{DpDph1}$};
\coordinate (group-13) at (2.600,-8.330);
\node[empty branch] (pending-13) at (2.600,-9.050) {\vzNoCovers\ no covers};
\coordinate (group-14) at (5.600,-8.330);
\node[contracted cover] (cover-14-0) at (5.600,-9.050) {\vzContracted\\[-1pt]$\mathit{DpDphh3}$};
\coordinate (group-15) at (-1.500,-10.230);
\node[cover] (cover-15-0) at (-5.900,-10.950) {\vzInconsistent};
\node[cover] (cover-15-1) at (-3.700,-10.950) {\vzInconsistent};
\node[cover] (cover-15-2) at (-1.500,-10.950) {\vzInconsistent};
\node[cover] (cover-15-3) at (0.700,-10.950) {\vzInconsistent};
\node[cover] (cover-15-4) at (2.900,-10.950) {\vzInconsistent};
\coordinate (group-16) at (5.900,-10.230);
\node[cover] (cover-16-0) at (5.900,-10.950) {\vzInconsistent};
\node[input] (input) at (0.000,0) {$P$};
\draw[stem] (input) -- (group-0);
\draw[branch] (group-0) -- (cover-0-0.north);
\draw[branch] (group-0) -- (cover-0-1.north);
\draw[branch] (group-0) -- (cover-0-2.north);
\draw[branch] (group-0) -- (cover-0-3.north);
\draw[stem] (cover-0-3.south) -- node[edge label,sloped,above] {$= 0$} (group-1);
\draw[branch] (group-1) -- (cover-1-0.north);
\draw[branch] (group-1) -- (cover-1-1.north);
\draw[stem] (cover-0-3.south) -- node[edge label,sloped,above] {$\ne 0$} (group-8);
\draw[branch] (group-8) -- (cover-8-0.north);
\draw[stem] (cover-1-1.south) -- node[edge label,sloped,above] {$= 0$} (group-2);
\draw[branch] (group-2) -- (cover-2-0.north);
\draw[branch] (group-2) -- (cover-2-1.north);
\draw[branch] (group-2) -- (cover-2-2.north);
\draw[stem] (cover-1-1.south) -- node[edge label,sloped,above] {$\ne 0$} (group-3);
\draw[branch] (group-3) -- (cover-3-0.north);
\draw[stem] (cover-3-0.south) -- node[edge label,sloped,above] {$= 0$} (group-4);
\draw[branch] (group-4) -- (pending-4.north);
\draw[stem] (cover-3-0.south) -- node[edge label,sloped,above] {$\ne 0$} (group-5);
\draw[branch] (group-5) -- (cover-5-0.north);
\draw[stem] (cover-5-0.south) -- node[edge label,sloped,above] {$= 0$} (group-6);
\draw[branch] (group-6) -- (cover-6-0.north);
\draw[branch] (group-6) -- (cover-6-1.north);
\draw[stem] (cover-5-0.south) -- node[edge label,sloped,above] {$\ne 0$} (group-7);
\draw[branch] (group-7) -- (cover-7-0.north);
\draw[stem] (cover-8-0.south) -- node[edge label,sloped,above] {$= 0$} (group-9);
\draw[branch] (group-9) -- (cover-9-0.north);
\draw[branch] (group-9) -- (cover-9-1.north);
\draw[branch] (group-9) -- (cover-9-2.north);
\draw[branch] (group-9) -- (cover-9-3.north);
\draw[stem] (cover-8-0.south) -- node[edge label,sloped,above] {$\ne 0$} (group-10);
\draw[branch] (group-10) -- (cover-10-0.north);
\draw[stem] (cover-10-0.south) -- node[edge label,sloped,above] {$= 0$} (group-11);
\draw[branch] (group-11) -- (cover-11-0.north);
\draw[branch] (group-11) -- (cover-11-1.north);
\draw[branch] (group-11) -- (cover-11-2.north);
\draw[branch] (group-11) -- (cover-11-3.north);
\draw[stem] (cover-10-0.south) -- node[edge label,sloped,above] {$\ne 0$} (group-12);
\draw[branch] (group-12) -- (cover-12-0.north);
\draw[stem] (cover-12-0.south) -- node[edge label,sloped,above] {$= 0$} (group-13);
\draw[branch] (group-13) -- (pending-13.north);
\draw[stem] (cover-12-0.south) -- node[edge label,sloped,above] {$\ne 0$} (group-14);
\draw[branch] (group-14) -- (cover-14-0.north);
\draw[stem] (cover-14-0.south) -- node[edge label,sloped,above] {$= 0$} (group-15);
\draw[branch] (group-15) -- (cover-15-0.north);
\draw[branch] (group-15) -- (cover-15-1.north);
\draw[branch] (group-15) -- (cover-15-2.north);
\draw[branch] (group-15) -- (cover-15-3.north);
\draw[branch] (group-15) -- (cover-15-4.north);
\draw[stem] (cover-14-0.south) -- node[edge label,sloped,above] {$\ne 0$} (group-16);
\draw[branch] (group-16) -- (cover-16-0.north);
\node[draw=gray!55,fill=white,rounded corners=2pt,inner sep=3pt,anchor=north] at (0,-11.750) {\textbf{Hitting-set outcome:}\quad \vzInconsistent\ UNSAT\quad \vzConsistent\ SAT\quad \vzTimeout\ timed out\quad \vzContracted\ contracted\quad \vzNoCovers\ no covers};
\end{tikzpicture}
}
    \caption{Recursion tree for \texttt{systems/system.txt}.  Each cover is shown separately.}
    \label{fig:variable_zeroing_tree}
\end{figure}

\subsection{Choice of variable for splitting}

The choice of variable for splitting effects the efficiency of the algorithm.
Every split creates at least two new covers to check.
If a variable barely occurs in the current cover, both recursive calls may repeat most of the same work.
A variable that occurs in many generators is a more promising choice: setting it to zero removes all terms containing it at once, while in the nonzero branch the algorithm can use the $1 - xz$.

For each candidate variable $x$ that is still available on the current branch, define
\begin{align*}
    g_Q(x) &= \#\{q\in Q : x\text{ occurs in }q\},\\
    t_Q(x) &= \#\{(q,m) : q\in Q,\ m\text{ is a term of }q \deg_x(m) > 0\},\\
    e_Q(x) &= \sum_{q\in Q}\ \sum_{\substack{m\text{ a term of }q \deg_x(m) > 0}} \deg_x(m).
\end{align*}
Its score is the triple
\[
    h_Q(x)=\bigl(g_Q(x),t_Q(x),e_Q(x)\bigr),
\]
and scores are compared lexicographically.

\subsection{Theoretical analysis of variable zeroing}

% Showing soundness and completeness of the method with variable zeroing

% Inference of a timeout evaluation result from its children evaluations' results
% if all its children are unsatisfiable, then the evaluation is unsatisfiable
% otherwise the evaluation is a union of the satisfiable children evaluations' results
